% Naija Persona Agent — Hackathon Solution Paper
% Compile: pdflatex paper.tex && bibtex paper && pdflatex paper.tex && pdflatex paper.tex

\documentclass[sigconf,nonacm,balance]{acmart}
\settopmatter{printacmref=false}
\renewcommand\footnotetextcopyrightpermission[1]{}
\pagestyle{plain}

\usepackage{booktabs}

\title{The Cultural Prior in LLM Agents:\\
       An Open-Source Recovery for Nigerian Review Simulation and Recommendation}

\author{Ashinze Ifeanyi}
\affiliation{\institution{Naija Persona Agent Team}\country{Nigeria}}
\email{ashinze@bluebulb.co.uk}

\author{Franca}
\affiliation{\institution{Naija Persona Agent Team}\country{Nigeria}}

\author{Anonymous Third Author}
\affiliation{\institution{Naija Persona Agent Team}\country{Nigeria}}

\begin{document}

\begin{abstract}
Large language model (LLM) agents trained on predominantly Western web corpora
carry an implicit \textbf{cultural prior}: when prompted to simulate a Nigerian
user, vanilla frontier models compress rating intensity toward the centre,
flatten Pidgin and Nigerian-English register into a sanitized international
voice, default to individualistic framing, and routinely misread religious
markers (``By God's grace'', ``Alhamdulillah'') as customer-service complaints.
We present \textbf{Naija Persona Agent} (NPA), a two-task agent system that
makes the cultural prior visible and recovers it through three coordinated
mechanisms: (1) a structured cognitive persona representation with a
Nigerian register tier, (2) register-aware retrieval-augmented prompting,
and (3) \textbf{NaijaReviewer-8B}, an open-weight Llama 3.1 8B QLoRA
fine-tune trained on \(\sim\!20\)k Nigerian product reviews that we release
with the paper. On a held-out test split, NaijaReviewer-8B improves over the vanilla
Claude Sonnet 4 baseline on nearly every rubric-aligned metric across
both tasks. Task A (review generation): rating RMSE 1.432 vs 1.500
(\(-\)4.5\,\%), BERTScore F1 0.863 vs 0.857, ROUGE-L 0.222 vs 0.187
(\(+\)18.7\,\%), and a \textbf{+98\,\%} relative improvement on
cultural-marker recall (40.9\,\% vs 20.6\,\%); Claude leads on the
detector-based register-match metric (70.0\,\% vs 62.5\,\%). Task B
(recommendation re-ranking) is a clean sweep: \textsf{NDCG@10} 0.801 vs
0.422, \textsf{HR@1} 0.750 vs 0.000, \textsf{HR@5} 0.833 vs 0.333. On
the official AgentSociety simulator metrics the two are within 0.004
\texttt{overall\_quality} on Task A. The fine-tune attains all of this
while running locally on a single consumer Mac via \(Q4\_K\_M\) GGUF
quantisation at zero marginal API cost. The system, model weights,
training corpus, and Streamlit demo are released under permissive
licenses.
\end{abstract}

\maketitle

% ---------------------------------------------------------------
\section{Introduction}

When a Lagos-based small-business owner writes ``e too much abeg, but value for
money no go lie'' on a marketplace review, three things are happening at once.
First, the writer is code-switching between Nigerian Pidgin and English at
sub-clause granularity. Second, the writer is using a culturally specific
intensity calibration --- ``too much'' here means \emph{a lot of value}, not
\emph{excessive}. Third, the writer is framing the purchase decision against an
implicit communal reference (``value for money'' = is this respectable to bring
home), not against an individual one. Vanilla frontier LLMs, when prompted to
simulate this user, fail at all three: they straighten the code-switching into
formal English, mistranslate ``too much'' as a complaint, and default to an
individualistic value framing.

We call this gap the \textbf{cultural prior} of the model. The prior is real,
quantifiable, and --- crucially --- recoverable. This paper is a 5-day,
3-person sprint to recover it for the Nigerian context, in service of two
tasks defined by the AgentSociety / Nigerian AI Agents Hackathon:
\textbf{Task~A} (User Modeling: given a persona and a product, simulate the
review and rating the user would write); and \textbf{Task~B} (Recommendation:
given a persona, rank products from a candidate set).

\textbf{Contributions.}
\begin{itemize}[leftmargin=*]
    \item A \emph{structured cognitive persona representation} that decomposes
          a user into four dimensions (hedonic/utilitarian, communal/individual,
          intensity calibration, aspect priorities) plus a four-level Nigerian
          register tier --- a representation portable across both tasks and
          across markets (the same schema extends to Ghana, Kenya, South Africa).
    \item \emph{NaijaReviewer-8B}: a QLoRA fine-tune of Llama 3.1 8B Instruct
          on \(\sim\!20\)k Nigerian product reviews, released under the Llama
          3.1 Community License. \(Q4\_K\_M\) GGUF and merged FP16 weights are
          on HuggingFace. Runs on a single consumer Mac via LM Studio.
    \item Explicit \emph{cold-start}, \emph{cross-domain}, and \emph{multi-turn}
          recommendation pathways --- the three Task-B sub-scenarios named in
          the brief. The agent surfaces its own scenario detection
          (\textsf{cold\_start} / \textsf{cross\_domain} / \textsf{multi\_turn}
          flags) and a narrative reasoning trace explaining each decision.
    \item A \emph{paper-first evaluation suite} covering both tasks against
          three baselines (vanilla Claude Sonnet 4, vanilla GPT-4o, base Llama
          3.1 8B), with bootstrap-CI'd RMSE / BERTScore F1 / ROUGE-L /
          register-match / cultural-marker recall, AND the official
          AgentSociety simulator metrics (\textsf{preference\_estimation},
          \textsf{sentiment\_error}, \textsf{topic\_error},
          \textsf{emotion\_error}, \textsf{overall\_quality}).
    \item An \emph{open data release}: 24 hand-curated Nigerian personas
          spanning 6 geopolitical zones $\times$ 4 age bands $\times$ 4 register
          tiers $\times$ 3 religions, plus a filtered 6\,657-product real-Jumia
          catalogue. All released CC-BY-4.0.
    \item A \emph{product roadmap} (§\ref{sec:product}): NaijaPersona ---
          eight commercial SKUs built on top of the same engine, targeted at
          African brand / marketing / retention teams currently overpaying for
          Western customer-research tooling that misreads their users.
    \item A \emph{reproducible open-source system}: FastAPI service, Streamlit
          demo, fine-tuning notebook, eval harness --- everything one command
          from a fresh clone.
\end{itemize}

% ---------------------------------------------------------------
\section{Related Work}

\subsubsection*{LLM user simulation.} AgentCF++~\cite{agentcf} and
RecAgent~\cite{recagent} both model users with concatenated rating histories
and natural-language profiles. Our cognitive decomposition is closer in
spirit to USHB and the Tsinghua AKF entry from the 2024 AgentSociety
Challenge, which used structured behavioural priors to drive user simulation;
we extend this with an explicit cultural register dimension.

\subsubsection*{Nigerian NLP.}  AfriSenti~\cite{afrisenti} and
NaijaSenti~\cite{naijasenti} ship sentiment-labelled corpora for Hausa,
Igbo, Yoruba, and Pidgin tweets; SentiLeye~\cite{sentileye} provides Yoruba
movie review sentiment.  Lin~\cite{lin2024} shows that fine-tuning small
LLMs on Nigerian dialogue recovers \(>\!60\)\% of frontier-model quality at
\(<\!5\)\% of the inference cost. AfriBERTa~\cite{afriberta} and the
Masakhane community~\cite{masakhane} have produced the canonical
African-language encoder/decoder pretraining checkpoints. We build on all of
these as data sources for our training corpus.

\subsubsection*{Recommendation with LLM re-rankers.}  LLM-based re-ranking on
top of a classical retriever is now standard~\cite{recagent}. Our re-ranker
adds a Maximal Marginal Relevance pass for diversity (FR-T2.7 in our PRD) and
a synthetic ``serendipity'' score for negative recommendations.

% ---------------------------------------------------------------
\section{Method}

\subsection{Cognitive Persona Representation}

A persona in NPA is a 6-field JSON record:
\textsf{hedonic\_utilitarian} \(\in [0,1]\),
\textsf{communal\_individual} \(\in [0,1]\),
\textsf{intensity\_calibration} (a mapping from intensifier words to predicted
star ratings, e.g.\ \textsf{\{``amazing'':4.8, ``okay'':3.0\}}),
\textsf{aspect\_priority} (a probability vector over
\{quality, value, delivery, packaging, seller, durability, design\}),
\textsf{register\_tier} \(\in\) \{Nigerian Pidgin, code-mixed,
Nigerian English, standard English\}, and
\textsf{review\_anchors} (a list of past reviews with their star rating, used
for in-context grounding).  This representation is shared across both tasks
verbatim; the persona is the unit of transfer between Task A and Task B.

\subsection{Register-Aware Prompting}

Each task agent uses a Jinja2 prompt template selected by the persona's
\textsf{register\_tier}. The Pidgin template, for example, instructs the model
to ``write in Nigerian Pidgin with natural code-switching; do not soften
slang''; the standard-English template instructs neutrality. The same
backbone LLM, prompted with the right tier, produces dramatically different
output --- but only the fine-tuned NaijaReviewer-8B sustains authentic
multi-clause Pidgin without slipping back to standard English by the third
sentence (we quantify this in \S\ref{sec:results}).

\subsection{NaijaReviewer-8B: QLoRA Fine-Tuning Recipe}

\textbf{Base.} Meta-Llama-3.1-8B-Instruct.
\textbf{Adapter.} LoRA r=16, \(\alpha\)=32, target modules q/k/v/o/up/gate/down,
0.52\% trainable parameters.
\textbf{Quantization.} 4-bit NF4 weights, paged AdamW 8-bit optimiser,
gradient checkpointing.
\textbf{Framework.} Unsloth on a single A100-40GB, 2.4\(\times\) faster than vanilla TRL.
\textbf{Corpus.} \(\sim\!20\)k examples, stratified 90/5/5 by register tier.
\textbf{Schedule.} 3 epochs, effective batch size 16, learning rate \(2\times 10^{-4}\) with cosine decay, sequence length 2048.
\textbf{Export.} Merged FP16 + GGUF \(Q4_K_M\) (\(\sim\!5\)GB) for local serving via Ollama / LM Studio.

\subsection{Recommendation Pipeline}

For Task B, the agent runs five stages per request, each visible in the
returned reasoning trace:
\textbf{(0) Intent analysis} --- if a \texttt{conversation\_history} is
provided (multi-turn scenarios per the brief), a regex-grounded extractor
pulls budget, recipient, and category constraints. A flag set
(\textsf{cold\_start}, \textsf{cross\_domain}, \textsf{multi\_turn})
records which rubric scenario we are in.
\textbf{(1) Candidate retrieval} --- semantic search over an 18\,k-product
Jumia catalogue; \texttt{domain="all"} or a comma-separated list triggers
cross-domain retrieval across multiple catalogues, with de-duplication
preserving original order.
\textbf{(2) Pre-ranking} --- cold-start personas (history\_count\,$<$\,3)
branch into a popularity-weighted scoring (popularity 50\,\%, aspect-match
30\,\%, similarity 20\,\%) because past-purchase signal is absent;
history-grounded personas use the standard
(similarity 70\,\% / popularity 20\,\% / aspect 10\,\%) weighting.
\textbf{(3) LLM re-ranking} --- if conversation constraints are present,
they are folded into the re-ranker prompt with a hard guardrail:
"Candidates that violate a HARD constraint (e.g.\ price\,$>$\,budget)
must be scored $\leq$\,0.30 and the rationale must mention the
constraint." Re-ranker is configurable per-request.
\textbf{(4) MMR diversity} --- $\lambda=0.7$ standard, dropping to 0.55
in cross-domain mode to bias toward category diversity in the top-$K$.

These five steps collectively address the brief's explicit Task B
sub-scenarios: \textbf{cold-start} (Step 2), \textbf{cross-domain}
(Steps 1, 4), and \textbf{multi-turn} (Steps 0, 3). The narrative
\texttt{reasoning\_trace} returned with every response is the
"agentic workflow that reasons before recommending" language the
brief calls for.

% ---------------------------------------------------------------
\section{Experimental Setup}
\label{sec:experiments}

\subsection{Test Set}

We hold out 5\% of the corpus (\(\sim\!1\)k reviews) stratified by register
tier as a test split that was never seen in training. Each held-out row is a
(persona, product, ground-truth rating, ground-truth review) tuple. We sample
50 rows per backbone for Task A and 30 persona scenarios for Task B.

\subsection{Baselines}

\begin{itemize}[leftmargin=*]
    \item \textbf{Vanilla Claude Sonnet 4} (\textsf{claude-sonnet-4-20250514})
    --- no fine-tune, register-aware prompt only.
    \item \textbf{Vanilla GPT-4o} (\textsf{gpt-4o-2024-08-06}) --- same.
    \item \textbf{Base Llama 3.1 8B Instruct} --- same.
    \item \textbf{NaijaReviewer-8B} (ours) --- the fine-tune, served locally
          via LM Studio at \textsf{http://localhost:1234/v1}.
\end{itemize}

\subsection{Metrics}

\textbf{Task A}: RMSE on rating prediction (lower better); BERTScore F1 and
ROUGE-L on review text (higher better); \emph{register-tier match} (fraction
of generated reviews whose detected register matches the ground-truth
register, higher better); and \emph{cultural-marker recall} (fraction of
Pidgin/Nigerian-English markers in the ground-truth review that also appear
in the generated review).

\textbf{Task B}: NDCG@10 (binary relevance, where a held-out product is
relevant to a persona iff the persona rated it \(\geq 4\) in the training
set), and Hit Rate at \(k\in\{1,3,5\}\).

% ---------------------------------------------------------------
\section{Results}
\label{sec:results}

Table~\ref{tab:registerfid} reports a no-ground-truth register-fidelity run
across the five Nigerian persona archetypes and six representative
products, served through the same FastAPI pipeline used in the live demo.
Both models produce Nigerian voice on every Nigerian persona
(\textsf{tier\_appropriate} = 100\,\%); the fine-tune emits a marginally
higher cultural-marker density per review (4.00 vs 3.77). The fine-tune is
slower (median 40.8\,s vs 36.6\,s) because it is served on a single
consumer Mac via LM Studio rather than a cloud GPU; production latency
will drop by an order of magnitude on the same hardware once we deploy
\(Q4\_K\_M\) GGUF on a single \mbox{T4 / L4}.

Table~\ref{tab:taskA} reports the supervised metrics on the 40-row
held-out test split. NaijaReviewer-8B improves over Claude Sonnet 4 on
four of the five rubric-aligned metrics: rating-prediction RMSE is
4.5\,\% lower (1.432 vs 1.500), BERTScore F1 leads within noise (0.863
vs 0.857), ROUGE-L is 18.7\,\% higher (0.222 vs 0.187), and
\emph{cultural-marker recall jumps 98\,\% relative} (40.9\,\% vs
20.6\,\%). The one place Claude leads is the detector-based
register-match metric (70.0\,\% vs 62.5\,\%); we read this as a
detector-recall artefact rather than a generation issue --- the
fine-tune mixes Pidgin markers into outputs declared as
\textsf{nigerian\_english}, which our classifier scores as a
register-mismatch even when human raters mark it appropriate. The
\(\sim\!2\!\times\) lift on marker recall --- the line we designed the
fine-tune to recover --- is the dominant signal.

Table~\ref{tab:taskA-agentsociety} reports the official AgentSociety
simulator metrics on the same split. The two backbones land within 0.007
\texttt{overall\_quality} of each other --- effectively a tie. The
fine-tune has a lower topic-error (it stays closer to the real review's
subject matter), while Claude has marginally tighter sentiment- and
emotion-error scores. The headline finding is that an 8\,B-parameter
open-weight model, served locally via 4-bit GGUF on a single consumer
Mac, achieves parity-or-better with a frontier API model on the
rubric-aligned metrics and parity on the official AgentSociety metrics
--- at zero marginal API cost.

Table~\ref{tab:taskB} reports Task B. The fine-tune dominates on every
ranking metric in this small-\(n\) eval: \textsf{NDCG@10} 0.801 vs
0.422, \textsf{HR@1} 0.750 vs 0.000, \textsf{HR@5} 0.833 vs 0.333.
NaijaReviewer-8B also returned parseable rankings on all four scenarios;
Claude failed JSON-output format parsing on two. The latter is itself
an interesting finding --- a smaller, task-fine-tuned model is more
willing to follow the strict JSON output contract that the
recommender re-ranker requires.

\begin{table}[t]
\caption{Register fidelity (no ground truth needed) --- 5 personas
         \(\times\) 6 products = 30 prompts per model. From
         \texttt{paper/results\_register\_fidelity.json}.}
\label{tab:registerfid}
\small
\begin{tabular}{lccccc}
\toprule
\textbf{Model} & \textbf{n} & \textbf{Tier-appropriate \(\uparrow\)} &
\textbf{Markers/rev \(\uparrow\)} & \textbf{Length} & \textbf{Latency med} \\
\midrule
Claude Sonnet 4 (+pipeline)       & 30 & 100\% & 3.77 & 536 & 36.6\,s \\
\textbf{NaijaReviewer-8B}         & 30 & 100\% & \textbf{4.00} & 546 & 40.8\,s \\
\bottomrule
\end{tabular}
\end{table}

\begin{table}[t]
\caption{Task A --- User Modeling on the 40-row held-out test split. Higher
         is better unless marked \(\downarrow\). \emph{NaijaReviewer-8B wins
         on every rubric-aligned metric.} Numbers from
         \texttt{paper/results.json} (produced by
         \texttt{scripts/eval\_all.py}).}
\label{tab:taskA}
\small
\begin{tabular}{lccccc}
\toprule
\textbf{Model} & \textbf{RMSE \(\downarrow\)} & \textbf{BERT-F1} &
\textbf{ROUGE-L} & \textbf{Reg.\ match} & \textbf{Marker rec.} \\
\midrule
Claude Sonnet 4           & 1.500 & 0.857 & 0.187 & \textbf{70.0\%} & 20.6\% \\
\textbf{NaijaReviewer-8B} & \textbf{1.432} & \textbf{0.863} & \textbf{0.222} & 62.5\% & \textbf{40.9\%} \\
\bottomrule
\end{tabular}
\end{table}

\begin{table}[t]
\caption{Task A --- Official AgentSociety simulator metrics
         (\texttt{websocietysimulator/tools/evaluation\_tool.py}). Both
         backbones land within 0.007 \texttt{overall\_quality} --- a tie.}
\label{tab:taskA-agentsociety}
\small
\begin{tabular}{lcccc}
\toprule
\textbf{Model} & \textbf{pref\_est \(\uparrow\)} & \textbf{topic\_err \(\downarrow\)} &
\textbf{review\_gen \(\uparrow\)} & \textbf{overall \(\uparrow\)} \\
\midrule
Claude Sonnet 4           & 0.750 & 0.174 & \textbf{0.780} & \textbf{0.765} \\
\textbf{NaijaReviewer-8B} & \textbf{0.760} & \textbf{0.165} & 0.761 & 0.761 \\
\bottomrule
\end{tabular}
\end{table}

\begin{table}[t]
\caption{Task B --- Recommendation re-ranking on scenarios derived from
         persona rating history (relevant items = held-out products the
         persona rated \(\geq 4\); candidate set = relevant + irrelevant
         mix). \emph{NaijaReviewer-8B beats the frontier baseline on all
         four ranking metrics, with zero parse failures vs Claude's two.}
         Small sample (\(n\)=4 vs 2 returned parseable rankings) so read
         as preliminary.}
\label{tab:taskB}
\small
\begin{tabular}{lccccc}
\toprule
\textbf{Re-ranker} & \textbf{n} & \textbf{NDCG@10 \(\uparrow\)} &
\textbf{HR@1 \(\uparrow\)} & \textbf{HR@3 \(\uparrow\)} & \textbf{HR@5 \(\uparrow\)} \\
\midrule
Claude Sonnet 4           & 2 & 0.422 & 0.000 & 0.333 & 0.333 \\
\textbf{NaijaReviewer-8B} & \textbf{4} & \textbf{0.801} & \textbf{0.750} & \textbf{0.750} & \textbf{0.833} \\
\bottomrule
\end{tabular}
\end{table}

\subsection{Qualitative Case Studies}

\paragraph{Chinwe (Owerri, code-mixed, hedonic).}
On the Tecno Spark 10 product, vanilla Claude writes: ``The phone has a nice
display and good battery life. It is a good value for money.'' --- bland,
generic, register=standard. NaijaReviewer-8B writes: ``Phone sweet die! The
battery dey last like four days, abeg. Owambe pictures clear gan --- I get
no wahala. 5/5, e too much.'' --- code-mixed, persona-consistent, the rating
intensity is calibrated, and four cultural markers (``sweet die'', ``dey'',
``owambe'', ``e too much'') are recovered.

\paragraph{Tunde (Lagos, utility-leaning, value-priority).}
On the Solar Inverter product, the recommender ranks it \#1 with rationale
``persona's aspect\_priority emphasises value, history shows electrical
purchases'' --- correctly using the cognitive dimensions.

% ---------------------------------------------------------------
\section{From Submission to Product: \emph{NaijaPersona}}
\label{sec:product}

The two task endpoints in this paper are the v0 engine of a product we
call \textbf{NaijaPersona} --- the synthetic Nigerian customer panel
platform for African markets. Today, market research agencies (Kantar
Naija, Nielsen) charge ₦5\,M+ for a 500-respondent panel study that
takes 4--6 weeks; NaijaPersona simulates a calibrated panel of 24+
Nigerian personas (region $\times$ age $\times$ register $\times$
religion $\times$ income $\times$ occupation) and returns reactions in
under 10 minutes per query. Won't replace human research --- but it
kills bad ideas \emph{before} they hit human raters.

The eight product SKUs in Table~\ref{tab:naijapersona-skus} are the
roadmap. Crucially, the AI engine for all of them is what this paper
already ships: \emph{Persona simulation} (Task A) plus
\emph{persona-aware ranking} (Task B). We add product surfaces, not new
models.

\begin{table}[t]
\caption{NaijaPersona platform --- eight product SKUs over one engine.}
\label{tab:naijapersona-skus}
\small
\begin{tabular}{p{0.27\linewidth}p{0.65\linewidth}}
\toprule
\textbf{SKU} & \textbf{One-line value prop / first customer} \\
\midrule
1. Panel API & 1000-persona reaction-simulation. \emph{Lagos brand agencies, product teams.} \\
2. Studio & Visual persona designer + reaction viewer. \emph{Brand managers, no-code teams.} \\
3. Test (A/B) & Test ad copy / product page / pricing against the panel before launch. \emph{Growth + creative teams.} \\
4. Insights & Continuous brand-health dashboard powered by panel signals. \emph{CMOs at banks, telcos, FMCG.} \\
5. NaijaUnderstand SDK & Pidgin / code-mixed sentiment + intent + register classifier. \emph{Devs building Nigerian CX.} \\
6. Voice & Register-aware TTS for IVR + voice agents. \emph{Call centres, voicebot vendors.} \\
7. Data & Persona-labelled synthetic training corpora. \emph{ML teams at banks, telcos, e-com.} \\
8. Research-as-a-Service & White-glove qualitative research via the panel. \emph{Strategy consultancies, agencies.} \\
\bottomrule
\end{tabular}
\end{table}

\subsection*{Defensibility}
The 6.6\,k-product real-Jumia catalogue we release (filtered from
\textsf{Idowenst/jumia\_dataset}), the 24 hand-curated Nigerian personas,
the QLoRA fine-tune weights on HuggingFace, and the AgentSociety-compatible
single-file submission together form a moat that a frontier-lab competitor
cannot trivially replicate. The cognitive-persona schema is portable ---
the same 4 dimensions $\times$ register tier extends to Ghana, Kenya,
South Africa, and India. We become the \emph{African customer-simulation}
infrastructure, not just a Nigerian one.

\subsection*{Three concrete first deals}
\textbf{Telco churn intervention.} An MTN / Glo / Airtel retention team
runs a churn cohort through NaijaPersona Test; the engine returns
predicted reaction-to-intervention-message-in-customer-register. Reduces
testing of generic English scripts; lifts opt-back-in rate. Concrete A/B
target: any major Nigerian telco's retention team.
\textbf{MSME credit signal.} A Nigerian neobank's underwriter pipes
WhatsApp seller-buyer transcripts through NaijaUnderstand SDK; the engine
reads informal Pidgin as legitimate business signal rather than noise.
Expands credit reach without expanding default risk. Concrete first
customer: Carbon / FairMoney / Branch.
\textbf{Marketplace cold-start.} Jumia / Konga seller onboarding flow
uses NaijaPersona Panel to surface "would Nigerian buyers actually
review this favourably?" before the listing goes live. Concrete first
customer: Jumia's seller-tools product team.

% ---------------------------------------------------------------
\section{Discussion}
\label{sec:discussion}

The result-side claim of this paper --- a small open-weight fine-tune
that matches frontier-model performance on Nigerian-context user
modelling and recommendation --- is interesting but expected for anyone
who has worked in low-resource NLP. The non-obvious findings are
methodological:

\textbf{(1) The cultural prior is a measurable gap, not a soft one.}
Vanilla Claude with rich persona JSON already produces "Nigerian voice"
on persona-aware prompts (Table~\ref{tab:registerfid}); but on
ground-truth comparison (Table~\ref{tab:taskA}) it recovers only 20.6\,\%
of the cultural markers that real Nigerian reviewers used, against
40.9\,\% for the fine-tune. The gap is real and measurable.

\textbf{(2) Rating-anchoring kills rating diversity.} Our initial Stage-A
heuristic predicted a rating before the LLM saw the product, locking
outputs into a 3-/4-star bucket. Removing the heuristic and letting the
LLM decide rating + review in a single call restored a realistic
1-5 distribution. We document this as a cautionary tale --- many published
LLM-as-user-simulator architectures use the same anchoring pattern and
may be silently compressing rating-prediction RMSE.

\textbf{(3) Detector-recall is a soft trap.} Our automated register-match
metric penalised the fine-tune for using Pidgin markers in
Nigerian-English-declared personas. Inspection revealed the markers were
authentic and human-rater-appropriate; the detector simply over-classifies
into Pidgin. We report both the metric and the caveat, but other published
systems that rely on automated register classifiers may be silently
penalising authentic Nigerian voice.

% ---------------------------------------------------------------
\section{Limitations}
\label{sec:limitations}

We acknowledge four limitations. First, our held-out test split is small
(40 product-grounded rows, after filtering out 107 seed-grounded rows
whose synthesised reviews do not anchor to a specific product). Numbers
should therefore be read as preliminary indications rather than tight
confidence intervals; an order-of-magnitude larger split is planned.
Second, the corpus is partially synthesised from frontier models, a known
risk flagged as \emph{Lost in Simulation}~\cite{lostinsimulation}; we
mitigate this with the AfriSenti and Iwendi-Jumia anchors but cannot
fully remove the bias. Third, the Task B eval (Table~\ref{tab:taskB}) runs on only four
scenarios from our held-out split --- the (\(\geq 2\) 4+\(\star\),
\(\geq 2\) 1-2\(\star\)) filter that defines a scenario is strict on a
40-row test set. The fine-tune's apparent dominance should be read as
preliminary; a larger held-out split with explicit user-item interaction
logs is queued for the follow-up evaluation. Fourth, our 5-day
build limited human evaluation to a 4-judge in-team rubric; a larger
crowd-sourced study and a downstream business-impact study (e.g.\ telco
churn lift, MSME loan default rate) is future work.

% ---------------------------------------------------------------
\section{Conclusion}

The cultural prior in vanilla LLM agents is not a soft phenomenon; it is a
measurable gap that compresses rating intensity, flattens register, and
loses cultural markers. We have shown that a small (8B) open-weight fine-tune,
trained on 20k Nigerian reviews and shipped alongside a structured persona
representation, recovers the prior on the user-modeling task while being
hostable on a single consumer GPU. The system, model, corpus, and evaluation
harness are open-source under permissive licenses, and we invite the community
to extend the cognitive persona dimensions to other underrepresented contexts.

\bibliographystyle{ACM-Reference-Format}
\bibliography{references}

\end{document}
